\chapter{The Physics}
\label{chap:the-physics}

% EPISODE COVERAGE: Episodes 9, 10, and 11 -- UniverseTensor, REAL, Variation,
% LOCAL, SpaceTimePath, CalculusProcess, UNIVERSAL, YarnTheory, HeartbeatProcess,
% LOGICAL, ElaborationProcess, HALTED
% KEY CONCEPT: Spacetime is the shape of the admissibility structure.
% Time is a ledger quantity, not a background.

\section{Three Derivatives}
\label{se:three-derivatives}

% Variation:
%   inductive Variation where
%     | newton: Variation
%     | gateaux: Variation
%     | frechet: Variation
%
% Newton: limit of secant slopes. The classical derivative.
% Gateaux: directional derivative in an infinite-dimensional function space.
%   Used in the calculus of variations.
% Frechet: the derivative as a bounded linear operator. The strongest notion.
%
% BigRedDogProcess holds a Variation and a transmute function.
% LOCAL: adds theory, delta, and experience.
%
% Sidebar: Why three derivatives? The three correspond to the three strategies
% in the instrument: Newton search (Cauchy sequences), Galerkin projection
% (finite element), and bisection (Cantor construction).

\begin{gphenom}{LiDAR}
\label{ph:lidar}
\GPhObs A LiDAR system measures distance by timing the return of a laser pulse.
Two LiDAR systems moving relative to each other will disagree on the timing
of the same pulse -- not because either is malfunctioning, but because each
has accumulated a different number of laser-certified distinctions.
\GPhCon Time dilation is not a physical distortion of time but a bookkeeping
disagreement between two ledgers. Each ledger is correct. They are uncorrelant
with respect to the timing event.
\GPhSeq The compiler that elaborates a \lean{SpaceTimePath.einstein} term is
keeping the same books as the LiDAR system. The elapsed heartbeats are the
certified distinctions. Two compilers running at different speeds are not
wrong -- they have different ledgers.
\end{gphenom}

\section{The Universe Tensor}
\label{se:universe-tensor}

% UniverseTensor:
%   structure UniverseTensor ... where
%     frame_of_reference: ...
%     reality: Truth
%     observe?: ...
%
% The accumulated linear encoding of all admissible refinements.
% Holds a Truth and a frame_of_reference.
% REAL wraps this with metaphysical?.
%
% LOCAL describes behavior in a neighborhood. REAL describes behavior globally.
% Together they provide the local-global structure of differential geometry.
%
% Sidebar: What is a tensor? A machine for measuring things that look different
% from different viewpoints but describe the same underlying reality.

\begin{gphenom}{Pound-Rebka}
\label{ph:pound-rebka}
\GPhObs Pound and Rebka (1959) sent gamma rays up a 22.5-meter shaft at Harvard
and measured the gravitational redshift. The photons arrived blueshifted
going down and redshifted going up by exactly the amount general relativity
predicts: $\Delta f / f = gh/c^2$.
\GPhCon The clock at the bottom of the tower accumulates distinctions more
slowly than the clock at the top. Both clocks are correct ledgers of what
they have witnessed. The gravitational potential is a ledger depth.
\GPhSeq The deeper you are in a gravitational well, the more slowly new
admissible events accumulate. Time dilation is the compiler running at
different heartbeat rates at different depths in the \lean{UniverseTensor}.
The tower has a gravitational gradient.
\end{gphenom}

\section{Spacetime as Data}
\label{se:spacetime-as-data}

% Episode 10: SpaceTimePath.
%
%   inductive SpaceTimePath where
%     | einstein: SpaceTimePath
%     | white_hole: SpaceTimePath
%     | blackhole: SpaceTimePath
%     | geodesic: SpaceTimePath
%
% einstein: a worldline in general relativity.
% white_hole: a time-reversed black hole. Matter and energy emerge.
% blackhole: information falls in and cannot return.
% geodesic: the path of least action, the straightest line in curved spacetime.
%
% CalculusProcess holds a derivative (a Variation) and a function (a SpaceTimePath).
% UNIVERSAL: the compiler, source_program, compiled_program, lake_build.
% The device describing itself. The tower is now aware of the tool that built it.
%
% Sidebar: What is a geodesic? In curved space, the straightest possible path.
% Light follows geodesics. Planets follow geodesics.

\begin{gphenom}{Hubble-Lemaitre}
\label{ph:hubble-lemaitre}
\GPhObs Lemaitre derived the expansion of the universe in 1927; Hubble confirmed
it observationally in 1929. Distant galaxies recede at a rate proportional to
their distance. The Hubble tension is the unresolved disagreement between
early-universe ($H_0 \approx 67$ km/s/Mpc) and late-universe
($H_0 \approx 73$ km/s/Mpc) measurements of this rate.
\GPhCon The expansion of the universe is the ledger growing. The rate of expansion
is the rate at which new admissible events can be recorded.
Early-universe and late-universe instruments have accumulated different ledgers
and are measuring the same rate from different positions in the record.
\GPhSeq The Hubble tension is structural, not experimental error. Two instruments
at different ledger depths disagree on the rate of ledger growth for the same
reason that two clocks at different gravitational depths disagree on elapsed time.
The disagreement is a measurement of the ledger's own curvature.
\end{gphenom}

\section{The Stokes Integral}
\label{se:stokes-integral}

% YarnTheory:
%   inductive YarnTheory where
%     | stokes: YarnTheory
%     | fibers: YarnTheory
%     | fabric: YarnTheory
%
% Stokes' theorem: the integral of a derivative over a region equals the
% integral of the original function over the boundary. Green's theorem,
% divergence theorem, and classical Stokes are all special cases.
%
% HeartbeatProcess: bullshit_meter, accumulated_bullshit: YarnTheory, weave?
% The elaboration overhead is modeled as a Stokes integral over the fabric
% of prior commitments. The bullshit meter is doing algebraic topology.

\section{The Halting of Physics}
\label{se:halting-of-physics}

% Episode 11: LOGICAL, ElaborationProcess, HALTED.
%
% HALTED returns True only for ComputerProgram.boolean -- the terminal decision.
% Not for load. Not for transform. Only boolean.
%
% A physical process terminates when it reaches a state where the outcome
% is a boolean. A geodesic terminates at a black hole or at a measurement.
% Load and transform are still in motion. Boolean has stopped.
%
% This closes the loop from Chapter 3: the halting problem is the physics of
% observation. A measurement is a process that halts. An unobserved system
% is a process that has not halted yet.

\begin{gphenom}{Photoelectric}
\label{ph:photoelectric}
\GPhObs Einstein showed in 1905 that the photoelectric effect depends on
frequency, not intensity: light below a threshold frequency cannot eject
electrons regardless of how bright it is, while light above the threshold
always can, even at very low intensity.
\GPhCon The admissibility cost of an event is set by frequency -- the possibility
of a distinction -- not by intensity -- the rate of attempts.
No number of sub-threshold attempts will force the compiler to accept a term.
\GPhSeq \lean{HALTED} returns \lean{True} only for \lean{ComputerProgram.boolean}.
The photoelectric threshold is the point at which the event crosses from
\lean{load} or \lean{transform} into \lean{boolean}: admissible, terminal,
witnessed. Below the threshold, the elaboration never halts.
\end{gphenom}

% End of chapter: the tower has agreed to hold spacetime paths, geodesics,
% white holes and black holes, Stokes integrals, compiler self-awareness,
% and the identification of physical halting with measurement.

